%!TEX program = uplatex
\RequirePackage{plautopatch}
\documentclass[uplatex,dvipdfmx]{jlreq}
\usepackage{amsmath,amssymb}

\pagestyle{empty}
\begin{document}

\noindent
{\large\bfseries トポロジーから見たトーリック多様体論}
\hfill
{\large\bfseries 枡田　幹也（大阪市立大学大学院理学研究科・理学部）}

\bigskip

「トーリック多様体」（toric variety）という代数幾何学の対象と
「扇」（fan）という組み合わせ論の対象の間に、１対１対応がある。
これにより、組み合わせ論の問題を代数幾何学の問題に翻訳し、
代数幾何学の定理（例えば Riemann--Roch、強 Lefschetz）を用いて解くということがしばしばあった。
R.~Stanley による $g$-予想の解決は、その最たるものであろう。
異なる数学分野を結ぶこのような架け橋に、私は魅了される。

本講演では、トーリック多様体論をトポロジーの立場から眺め、
トーリック多様体論のトポロジー版を考える。
つまり、トポロジーと組み合わせ論の間に架け橋を構築することを試みる。

トポロジーの観点からすると、対象をトーリック多様体に限る必要はない。
「ユニタリトーリック多様体」と呼ばれるもっと広い幾何学的対象を取り扱うのが自然である。
その際現れてくる組み合わせ論の対象は、扇が幾重にも重なった「多重扇」（multi-fan）である。
この多重扇の重なりの次数は、Todd 種数に対応する。
よく知られているように、トーリック多様体の Todd 種数は１である。
（我々の立場からすると）それ故、トーリック多様体の扇は重なりのないものとなると解釈できる。

我々の議論では、同変コホモロジーが中心的な役割を演ずる。
これは1960年頃 A.~Borel によって導入され、群作用の情報をよく含んでいることが知られている。
また、同変 $K$ 群などより取り扱いやすいという利点がある。
しかし、一部の専門家を除いては、同変コホモロジーはあまりポピュラーではないようである。
１回目の講演では、同変コホモロジーが（ユニタリ）トーリック多様体と相性が良いことを観る。
我々の多重扇は、同変コホモロジーを用いて定義される。
２回目の講演では、Riemann--Roch 指数の重複度公式、Pick 公式の一般化等を考える。

以上の書き方では、我々の議論がトーリック多様体論を含んだ形で展開されているような
印象を与えるが、そうではない。以下の２つの問題点がある。
\begin{itemize}
  \item 我々は、コンパクトで滑らかな多様体しか扱っていない。
        一方、トーリック多様体論では、コンパクトでないものや特異点を許すものも扱っている。
  \item トーリック多様体論は、トーリック多様体と扇との間の「両方通行の架け橋」である。
        しかし、（今の所）我々の理論は、ユニタリトーリック多様体に多重扇を対応させることが
        出来るというだけの「一方通行の架け橋」である。
        勝手に多重扇を与えたとき、それに対応するユニタリトーリック多様体が
        （ある意味で）唯一つだけあるということは、判っていない。
\end{itemize}
講演の最後に、これらの問題点、今後の研究課題に触れたいと思っている。

\end{document}